\section{完整 DMAIC 可用性提升设计}

\begin{frame}{DMAIC 总览：先修测量，再调策略，最后固化控制}
\centering
\resizebox{0.96\linewidth}{!}{%
\begin{tikzpicture}[
  phase/.style={draw,rounded corners=3pt,minimum width=2.35cm,minimum height=1.15cm,align=center,font=\small\bfseries,thick},
  a/.style={-{Latex[length=2mm]},thick,Navy},node distance=0.35cm
]
\node[phase,fill=SoftBlue] (d) {D - Define\\覆盖合同/用户价值};
\node[phase,fill=SoftGreen,right=of d] (m) {M - Measure\\MSA/基线/数据计划};
\node[phase,fill=SoftAmber,right=of m] (a1) {A - Analyze\\根因/假设/优先级};
\node[phase,fill=SoftRed,right=of a1] (i) {I - Improve\\实验/实现/试点};
\node[phase,fill=Paper,right=of i] (c) {C - Control\\SPC/回滚/治理};
\draw[a] (d)--(m);\draw[a] (m)--(a1);\draw[a] (a1)--(i);\draw[a] (i)--(c);
\end{tikzpicture}
}

\vspace{1em}
\begin{tabularx}{\textwidth}{p{0.14\textwidth} X X}
\toprule
Tollgate & 必须交付 & 放行条件 \\
\midrule
D & charter、SIPOC、VOC-CTQ、范围与风险 & “可用 corpus”定义无歧义 \\
M & 冻结 benchmark、操作定义、MSA、现状能力 & 标签/事件/主题测量达到门槛 \\
A & Pareto、因果图、验证过的根因 & 至少一条证据支持每个优先根因 \\
I & DOE/A-B、pilot、FMEA、变更包 & 主要 CTQ 改善且无安全/成本回退 \\
C & control plan、看板、告警、回滚与审计 & 连续 4 周/多个 campaign 稳定 \\
\bottomrule
\end{tabularx}
\end{frame}

\begin{frame}{Define 1：项目章程与问题陈述}
\begin{columns}[T,onlytextwidth]
\column{0.49\textwidth}
\begin{block}{Problem statement}
当前系统能稳定采集与持久化，但 latest full 仅从 NVD 获得 28 条、15 个 Core quota gap 未满足；全库 67.4\% uncertain、75.0\% 来自 arXiv。现有 full 终态不能证明 corpus 的质量、证据结构或可恢复性。
\end{block}
\begin{block}{Business case}
降低人工检索与复核成本；更快发现 LLM/agent 风险；给 item-KG/global KG 提供可追溯、结构平衡的输入。
\end{block}
\column{0.47\textwidth}
\begin{block}{Goal statement（建议）}
在冻结预算与源快照上，将 quality-validated weighted coverage 提升到 $\ge0.90$，无 critical gap；人工 precision $\ge0.85$；uncertain $\le0.25$；100\% campaign 可恢复；单位有效证据成本不高于基线 1.25 倍。
\end{block}
\begin{block}{Scope}
In：intel\_agent full/incremental、四源、relevance、Mongo/trace、item-KG handoff。Out：base KG 本体重构、生产 SOC 响应、全自动漏洞处置。
\end{block}
\end{columns}
\end{frame}

\begin{frame}{Define 2：VOC -> CTQ 与 SIPOC}
\scriptsize
\begin{tabularx}{\textwidth}{p{0.15\textwidth} p{0.23\textwidth} X p{0.22\textwidth}}
\toprule
Voice & 需要 & CTQ & 缺陷定义 \\
\midrule
安全分析师 & 少噪声、可追溯、可复核 & precision、provenance completeness & 无法证实来源或错误相关 \\
研究人员 & 主题与研究/漏洞证据都覆盖 & weighted evidence coverage、independence & critical cell 未满足 \\
运维 & 可控成本、可恢复、失败不静默 & cost/evidence、resume success、fallback visibility & 丢水位/静默降级/超预算 \\
KG 消费方 & DTO 完整、实体锚点、文本质量 & handoff completeness、ID density & 缺 raw/summary/source/ID \\
治理方 & 安全工具使用、版本可审计 & policy violation rate、trace completeness & 越权调用或 trace 缺失 \\
\bottomrule
\end{tabularx}

\vspace{0.5em}
\begin{block}{SIPOC}
\textbf{Suppliers}: NVD/CISA/OSV/arXiv、LLM provider、taxonomy owner
$\rightarrow$ \textbf{Inputs}: API records、goal、budget、checkpoint、playbook
$\rightarrow$ \textbf{Process}: plan-call-parse-relevance-gap-stop-persist
$\rightarrow$ \textbf{Outputs}: corpus、gap state、trace、item-KG queue
$\rightarrow$ \textbf{Customers}: analyst、researcher、KG、console/ops。
\end{block}
\end{frame}

\begin{frame}{Measure 1：指标操作定义 - 避免同名异义}
\scriptsize
\begin{tabularx}{\textwidth}{p{0.22\textwidth} X p{0.24\textwidth}}
\toprule
CTQ & 计算定义 & 数据源/频率 \\
\midrule
人工 precision & 双标注 relevant / 被系统判 relevant & 每版本分层抽样 \\
Recall proxy & 冻结 gold events 中被自然键或人工语义匹配召回的比例 & 每次 benchmark \\
Weighted coverage & $1-\sum w_gG^{cov}_g/\sum w_g$；另报 critical gap count & 每 run/每轮 \\
Independent evidence & canonical events / relevant items；异源佐证另报 & 每 run \\
Freshness compliance & 有效证据最新时间满足 channel SLA 的 cell 比例 & 每日 \\
Efficiency & $\Delta$ effective evidence / source API call；并报美元与秒 & 每轮/每 run \\
Handoff completeness & 必需 DTO 字段完整且可入 KG queue 的 item 比例 & 每 run \\
Resumability & 注入中断后无遗漏/无水位跳跃完成的 campaign 比例 & 每发布 \\
Safety & 越权 tool call、未审计 fallback、恶意内容导致违规动作的比例 & 每安全套件 \\
\bottomrule
\end{tabularx}
\end{frame}

\begin{frame}{Measure 2：MSA - relevance、topic 与 event 三套量具分别验证}
\scriptsize
\begin{columns}[T,onlytextwidth]
\column{0.49\textwidth}
\begin{block}{抽样设计}
\begin{itemize}
  \item $n=600$，按 4 sources × 初判标签 × 时间层分层
  \item 2 名安全标注员独立判断；冲突由第 3 人裁决
  \item 14 天后对 15\% 样本重复标注，测 repeatability
  \item 保留 title/summary/raw/source，不显示系统 score 以减少锚定
\end{itemize}
\end{block}
\column{0.47\textwidth}
\begin{block}{量具与门槛（建议）}
\begin{itemize}
  \item relevance：Cohen $\kappa\ge0.75$
  \item multi-label topic：macro F1 $\ge0.80$，Jaccard $\ge0.75$
  \item canonical event：pair agreement $\ge0.95$
  \item score calibration：ECE $\le0.05$，Brier score 报告 CI
  \item 若不达标：修订 rubric/anchor，再测，不进入 Analyze
\end{itemize}
\end{block}
\end{columns}

\begin{alertblock}{关键纪律}
uncertain 不计作健康覆盖。若业务允许软证据，必须以校准后的概率和置信区间传播，而不是把标签名称直接变成常数权重。
\end{alertblock}
\end{frame}

\begin{frame}{Measure 3：基线数据计划与分层}
\scriptsize
\begin{tabularx}{\textwidth}{p{0.18\textwidth} X X}
\toprule
数据集 & 设计 & 目的 \\
\midrule
Frozen-source replay & 保存源响应/时间戳/分页，30 个任务 × 5 seeds & 消除 live source 漂移，比较算法 \\
Live shadow & 同一天、同预算交错运行策略，不自动入主库 & 验证外部有效性与时延/限流 \\
Adversarial source & 注入 prompt、伪 CVE、重复事件、污染 metadata & 验证 hooks、parser、relevance、安全 \\
Interruption suite & 在分页、persist、SDK session、Mongo 处注入故障 & 验证 checkpoint/resume/fallback \\
Human gold & 600 MSA + 每任务事件/主题 gold 子集 & precision、recall proxy、calibration \\
KG handoff & 固定 100 item，含空 raw、同义实体、冲突 provenance & DTO 完整与下游质量 \\
\bottomrule
\end{tabularx}

\vspace{0.2em}
\begin{block}{分层原则}
所有指标按 source、topic rarity、channel、query intent、provider、fallback state、time window 分层；总均值只作摘要，不能掩盖长尾主题和单源失败。
\end{block}
\end{frame}

\begin{frame}{Analyze 1：当前证据支持的因果假设}
\small
\begin{tabularx}{\textwidth}{p{0.28\textwidth} X p{0.22\textwidth}}
\toprule
假设 & 现有证据 & 下一步验证 \\
\midrule
H-A1 full 目标函数错位 & full 条件在 controller 明确排除 adaptive；latest 有 15 quota gap 且 0 CorpusGap & 同预算 legacy vs campaign \\
H-A2 relevance 量具失真 & uncertain 67.4\%；LLM 判定仅 0.35\%；raw score 从 .45 起 & MSA + threshold/calibration \\
H-A3 吞吐 reward 引起 source skew & arXiv 占 75.0\%；高吞吐源天然赢 new-items/call & source-normalized reward 消融 \\
H-A4 stop 缺乏交付合同 & last action STOP 与 stop\_reason null；未满足 gap 仍停 & risk-bounded stop A/B \\
H-A5 schema/看板语义导致误判 & core\_gap\_count=0 与 15 legacy gaps 并存 & contract tests + dashboard audit \\
H-A6 playbook 学习信号稀疏 & 22 条中仅 1 verified、13 deprecated & replay with/without memory \\
\bottomrule
\end{tabularx}
\end{frame}

\begin{frame}{Analyze 2：Pareto + FMEA，而不是凭直觉排期}
\begin{columns}[T,onlytextwidth]
\column{0.47\textwidth}
\begin{block}{缺陷 Pareto 建议分类}
\begin{enumerate}
  \item false relevant / uncertain overflow
  \item critical cell uncovered
  \item source/channel imbalance
  \item premature stop / budget waste
  \item checkpoint or persistence gap
  \item missing provenance / KG handoff defect
  \item policy violation / silent fallback
\end{enumerate}
按对 CTQ 的损失而非事件次数排序。
\end{block}
\column{0.49\textwidth}
\begin{block}{FMEA 评分}
对 failure mode 评 Severity $S$、Occurrence $O$、Detection $D$，$RPN=SOD$；另强制将任何越权工具调用、checkpoint 跳跃、provenance 丢失设为 stop-ship，不允许被低频稀释。
\end{block}
\begin{alertblock}{因果验证}
使用冻结 replay 改一个因子；live shadow 做外部验证。相关性只说明线索，不等于根因。
\end{alertblock}
\end{columns}
\end{frame}

\begin{frame}{Improve 1：五个改进包与优先顺序}
\scriptsize
\begin{tabularx}{\textwidth}{p{0.11\textwidth} p{0.22\textwidth} X p{0.18\textwidth}}
\toprule
优先级 & 改进包 & 具体动作 & 成功门槛 \\
\midrule
P0 & Measurement contract & MSA、probability calibration、统一 stop/gap schema、data contract tests & 量具门槛全通过 \\
P1 & Full campaign & full 复用 CorpusGap/UCB/checkpoint；阶段化 seed-gap-confirm-discover-handoff；可 resume & coverage $\ge.90$，无 critical \\
P2 & Source-aware policy & applicability mask、cost/latency normalized contextual bandit、portfolio attribution & 同预算 evidence/call 提升 \\
P3 & Active relevance & uncertainty sampling、source/topic calibration、LLM escalation budget、hard negatives & uncertain $\le.25$，precision $\ge.85$ \\
P4 & KG closed loop & schema mapping、entity alignment、conflict/provenance、KG gap -> query & 图完整度/下游任务提升 \\
\bottomrule
\end{tabularx}

\vspace{0.6em}
\begin{alertblock}{发布原则}
每个改进包先 frozen replay，再 shadow，最后 5\% canary；任何 safety、precision 或 resumability stop-ship 指标回退即自动 rollback。
\end{alertblock}
\end{frame}

\begin{frame}{Improve 2：把 full 重构为可恢复的阶段化 Bootstrap Campaign}
\centering
\begin{tikzpicture}[
  n/.style={draw,rounded corners,minimum width=2.25cm,minimum height=0.9cm,align=center,font=\scriptsize\bfseries,thick},
  a/.style={-{Latex[length=1.7mm]},thick,Navy},node distance=0.35cm
]
\node[n,fill=SoftBlue] (seed) {1 Seed sweep\\适用源最小探测};
\node[n,fill=SoftGreen,right=of seed] (gap) {2 Gap fill\\高 P cell};
\node[n,fill=SoftAmber,right=of gap] (confirm) {3 Confirm\\CISA/异源佐证};
\node[n,fill=SoftRed,right=of confirm] (discover) {4 Discover\\Extended/residual};
\node[n,fill=Paper,right=of discover] (handoff) {5 Handoff\\质量清单/manifest};
\draw[a] (seed)--(gap);\draw[a] (gap)--(confirm);\draw[a] (confirm)--(discover);\draw[a] (discover)--(handoff);
\draw[a,dashed,bend left=30] (discover.north) to node[above,font=\tiny]{新 candidate 经审批进入 taxonomy backlog} (gap.north);
\end{tikzpicture}

\vspace{0.8em}
\begin{block}{Campaign stop contract}
只有同时满足：no critical gap；weighted coverage 达标；precision 下置信界达标；checkpoint/run gaps 清零或显式 waiver；预算/成本在界内；manifest 与 trace 完整，才可标记 complete。
\end{block}

\begin{block}{可恢复状态}
每个阶段、cell、candidate、source page、watermark、failure、model/prompt/taxonomy version 都持久化；重启从最后 confirmed checkpoint 继续，而非重新“full”。
\end{block}
\end{frame}

\begin{frame}{Control：Control Plan、SPC、告警与反应规则}
\scriptsize
\begin{tabularx}{\textwidth}{p{0.18\textwidth} p{0.18\textwidth} p{0.17\textwidth} X}
\toprule
指标 & 图/频率 & 触发条件 & 反应计划 \\
\midrule
人工 precision & p-chart / 每周样本 & LCL 低于 .85 或连续 2 点下行 & 停止自动入库，回滚阈值/模型，复核最近批次 \\
uncertain share & p-chart / 每 run & $> .25$ 或 source 单层异常 & 检查 embedding/LLM escalation、漂移样本 \\
weighted coverage & run chart / 每轮 & critical gap 两轮不降 & 重排 portfolio，打开专家 review \\
evidence/API call & I-MR / 每轮 & 低于基线 LCL & 冻结该 arm，检查 query/source 限流 \\
duplicate/noise & p-chart / 每源 & 超 UCL & 调整 canonicalization/query 或短路该源 \\
resume success & release gate & 任一水位跳跃/遗漏 & stop-ship，恢复旧版本并审计 checkpoint \\
policy violations & zero-tolerance & 任一越权/未审计 fallback & stop-ship、安全事件流程 \\
\bottomrule
\end{tabularx}

\vspace{0.5em}
\begin{block}{治理与可追溯}
taxonomy、prompt、provider/model、threshold、source client、schema 与 playbook 全部版本化；看板同时显示 spec limit 与统计 control limit；每次 waiver 要有 owner、到期日和证据。
\end{block}
\end{frame}

\begin{frame}{DMAIC 完成定义：不是上线一个功能，而是证明过程能力}
\small
\begin{columns}[T,onlytextwidth]
\column{0.48\textwidth}
\begin{exampleblock}{项目可以关闭的证据}
\begin{itemize}
  \item MSA 与 frozen benchmark 已发布
  \item 预注册同预算实验达到 CTQ，含 CI 与效应量
  \item adversarial/interruption 套件全通过
  \item campaign 从中断恢复无遗漏
  \item Control Plan 连续多个周期无 special cause
  \item owner、runbook、rollback、version manifest 完整
\end{itemize}
\end{exampleblock}
\column{0.48\textwidth}
\begin{alertblock}{不算完成的证据}
\begin{itemize}
  \item item 数增长
  \item 单次 live run 成功
  \item 单元测试通过但无人工 gold
  \item 平均 coverage 高但有 critical gap
  \item 模型理由“认为已经足够”
  \item 看板绿灯但字段口径未统一
\end{itemize}
\end{alertblock}
\end{columns}

\vspace{0.2em}
\begin{center}
\large\textcolor{Purple}{DMAIC 的终点是稳定的交付过程，不是一次漂亮的 demo。}
\end{center}
\end{frame}
